% ════════════════════════════════════════════════════════════════════════════
% content/chapters/chapter3.tex
% ════════════════════════════════════════════════════════════════════════════

\chapter{Hito 3: Modelado, experimentación y validación}

% ─────────────────────────────────────────────────────────────────────────────
\section{Introducción al Hito 3}

Tras haber comprendido la naturaleza fisicoquímica de los datos y su estructura multivariante en el Hito 2, el presente capítulo aborda la fase central del sistema VinOps: el \textbf{modelado predictivo}. El objetivo principal es diseñar, entrenar y validar un conjunto de algoritmos de \textit{Machine Learning} capaces de clasificar de forma robusta la variedad de los cultivares de vino.

Para garantizar la viabilidad del sistema en un entorno de producción real, este proceso no se limita a la búsqueda de la máxima exactitud, sino que sigue una metodología científica estricta: preprocesamiento fundamentado, diseño de experimentos controlados mediante ajuste de hiperparámetros (\textit{hyperparameter tuning}), evaluación mediante métricas robustas al desbalanceo y, finalmente, un análisis de interpretabilidad para validar las reglas subyacentes del modelo.
% ─────────────────────────────────────────────────────────────────────────────
\section{Diseño experimental y metodología de validación}

En esta fase se ha abordado el problema de clasificación multiclase sobre el \textit{Wine Recognition Dataset} mediante un diseño experimental orientado a maximizar la capacidad de generalización de los modelos. Para ello, se ha realizado una partición estratificada de los datos, destinando el 70\% de las observaciones al entrenamiento y el 30\% al test, preservando en ambos subconjuntos la proporción original de las tres clases de vino.

Sobre el conjunto de entrenamiento se ha aplicado una validación cruzada de 10 particiones (\textit{10-fold Cross-Validation}) durante el ajuste de los modelos. Esta estrategia permite evitar la necesidad de reservar un conjunto de validación independiente, algo especialmente conveniente en escenarios de \textit{small data}, donde resulta prioritario aprovechar al máximo la información disponible. En consecuencia, el conjunto de test queda completamente aislado y se utiliza exclusivamente para la evaluación final.

Siguiendo los criterios definidos en el Hito 2, la comparación entre modelos se ha basado principalmente en el \textbf{Macro F1-Score} y la \textbf{Balanced Accuracy}, complementadas por \textit{Accuracy}, \textit{Kappa} y el coeficiente \textit{Matthews Correlation Coefficient} (MCC). Este conjunto de métricas permite evaluar el rendimiento de forma robusta en un contexto multiclase con ligero desbalanceo.

% ─────────────────────────────────────────────────────────────────────────────
\section{Selección de algoritmos y ajuste de hiperparámetros}

La selecci\'on de algoritmos responde a una estrategia en dos niveles. Por
un lado, se incluyen los modelos presentes en los benchmarks de referencia
identificados en el Hito~2 (LDA, Regresi\'on Log\'istica Multinomial, SVM
radial, kNN y \'Arbol de Decisi\'on), cuya inclusi\'on permite contrastar
directamente los resultados obtenidos con los reportados en la literatura.
Por otro lado, el equipo propone \textbf{Random Forest como modelo principal
del sistema VinOps}, por las siguientes razones t\'ecnicas:

\begin{itemize}
    \item \textbf{Robustez ante multicolinealidad:} El an\'alisis de
    correlaciones del Hito~2 revel\'o un bloque fenólico altamente
    correlacionado (TotalPhenols, Flavanoids, OD280/OD315, $|r| \geq 0{,}70$).
    Random Forest es inherentemente robusto ante multicolinealidad gracias
    a la selecci\'on aleatoria de variables en cada divisi\'on (\texttt{mtry}),
    a diferencia de modelos lineales como la regresi\'on log\'istica que pueden
    verse afectados por inestabilidad num\'erica.

    \item \textbf{No requiere normalización:} A diferencia de SVM y kNN,
    Random Forest no es sensible a la escala de las variables. Dado que el
    dataset presenta escalas muy heterogéneas (Proline: 278--1680\,mg/L;
    Hue: 0{,}48--1{,}71\,u.a.), esta propiedad reduce el riesgo de introducir
    sesgos durante el preprocesado.

    \item \textbf{Interpretabilidad mediante importancia de variables:}
    El sistema VinOps necesita ofrecer al en\'ologo una explicaci\'on de 
    qu\'e variables han determinado la clasificaci\'on. Random Forest 
    proporciona de forma nativa una medida de importancia de
    variables basada en la reducci\'on del criterio de impureza (índice de Gini),
    directamente interpretable desde el dominio vitivin\'icola.

    \item \textbf{Resistencia al sobreajuste en \textit{small data}:}
    Al promediar sobre m\'ultiples \'arboles entrenados en subconjuntos
    aleatorios (\textit{bagging}), Random Forest reduce la varianza del
    estimador, mitigando el riesgo de sobreajuste que el reducido tama\~no
    muestral hace especialmente cr\'itico \cite{Breiman2001}.
\end{itemize}

El resto de modelos se incluyen como referencias experimentales que permiten
cuantificar la ganancia obtenida frente a alternativas m\'as simples y
contextualizar el rendimiento de Random Forest en el espectro de la
literatura.


Se han entrenado y evaluado siete modelos representativos, abarcando enfoques lineales, probabilísticos, basados en distancias y en árboles de decisión.

\begin{itemize}
    \item \textbf{LDA y Regresión Logística Multinomial:} Empleados como modelos paramétricos base por su solidez estadística y buena interpretabilidad.
    \item \textbf{Naive Bayes (NB):} Se ha optimizado el uso de estimaciones de densidad no paramétricas frente a distribuciones gaussianas normales (\textit{usekernel}).
    \item \textbf{K-Vecinos Más Cercanos (kNN):} Se exploró el número de vecinos óptimos ($k \in \{2, 3, 5, 8, 10\}$) utilizando distancia euclídea.
    \item \textbf{Máquinas de Vectores de Soporte (SVM Radial):} Dada la posible no linealidad de las fronteras químicas, se optimizó tanto el coste de penalización ($C \in \{0.1, 1, 10, 100\}$) como el parámetro del \textit{kernel} radial ($\sigma$).
    \item \textbf{Árboles (CART) y Random Forest (RF):} Modelos robustos ante la multicolinealidad detectada en el bloque fenólico. Para RF, se ajustó el número de variables seleccionadas aleatoriamente en cada división ($mtry \in \{2, 4, 6, 8, 10\}$).
\end{itemize}

El ajuste de hiperparámetros se ha realizado mediante Grid Search dentro del esquema de validación cruzada. Esta elección se justifica por el reducido tamaño del dataset y por el carácter acotado de los espacios de búsqueda definidos, lo que hace viable una exploración exhaustiva sin necesidad de recurrir a técnicas más complejas como \textit{Random Search} u optimización bayesiana. Además, la búsqueda en rejilla facilita la trazabilidad experimental y la interpretación metodológica de los resultados.

Las rejillas de b\'usqueda se dise\~naron espec\'ificamente para cada
modelo, seleccionando \'unicamente los hiperpar\'ametros con mayor impacto
conocido en el rendimiento y definiendo rangos acotados y justificados.
La Tabla~\ref{tab:grid} detalla, para cada modelo, el hiperpar\'ametro
optimizado, la justificaci\'on de su elecci\'on y el espacio de b\'usqueda
utilizado.

\begin{table}[H]
    \centering
    \caption{Justificaci\'on del espacio de b\'usqueda de hiperpar\'ametros
    por modelo}
    \label{tab:grid}
    \small
        \begin{tabularx}{\textwidth}{llXl}
        \toprule
        \textbf{Modelo} & \textbf{Hiperpar\'ametro} & \textbf{Justificaci\'on} &
        \textbf{Espacio de b\'usqueda} \\
        \midrule
        Naive Bayes
            & \texttt{usekernel}
            & Controla si la densidad condicional se estima con una Gaussiana
            (FALSE) o con estimaci\'on no param\'etrica por kernel (TRUE). Se
            explora porque las distribuciones de MalicAcid y Proline presentan
            asimetr\'ia notable, lo que puede violar la hip\'otesis gaussiana.
            & \{TRUE, FALSE\} \\
        \addlinespace
        kNN
            & $k$ (n.\textsuperscript{o} de vecinos)
            & Controla el equilibrio sesgo-varianza: valores peque\~nos de $k$
            producen fronteras muy locales (alta varianza); valores grandes
            suavizan la frontera (mayor sesgo). El rango {2--10} cubre desde
            fronteras muy precisas hasta una media local robusta, adecuado para
            178 muestras.
            & $k \in \{2, 3, 5, 8, 10\}$ \\
        \addlinespace
        SVM radial
            & $C$ (coste de penalizaci\'on)
            & Regula el n\'umero de vectores de soporte: $C$ bajo tolera m\'as
            errores de clasificaci\'on (mayor margen, mayor regularizaci\'on);
            $C$ alto busca la separaci\'on perfecta (riesgo de sobreajuste).
            El rango logar\'itmico cubre desde modelos muy regularizados hasta
            modelos muy ajustados.
            & $C \in \{0{,}1,\, 1,\, 10,\, 100\}$ \\
            & $\sigma$ (ancho del kernel RBF)
            & Controla el radio de influencia de cada muestra de entrenamiento.
            Un $\sigma$ bajo produce fronteras muy locales; uno alto las suaviza.
            Se utiliza estimaci\'on autom\'atica (sigest) como punto de partida
            y se refina manualmente.
            & Estimado autom\'aticamente \\
        \addlinespace
        \'Arbol (CART)
            & \texttt{cp} (par\'ametro de complejidad)
            & Controla la poda del \'arbol: un \texttt{cp} alto produce \'arboles
            m\'as simples (menor sobreajuste); uno bajo permite \'arboles m\'as
            profundos. El rango explorado abarca desde \'arboles muy podados
            hasta \'arboles de profundidad media, adecuado para el tama\~no del
            dataset.
            & $\{0{,}01,\, 0{,}05,\, 0{,}07,\, 0{,}10\}$ \\
        \addlinespace
        Random Forest
            & \texttt{mtry} (variables por divisi\'on)
            & N\'umero de variables seleccionadas aleatoriamente en cada
            divisi\'on del \'arbol. Valores bajos aumentan la diversidad entre
            \'arboles (menor correlaci\'on, menor varianza del ensamblado);
            valores altos se aproximan a un \'arbol \'unico. La regla emp\'irica
            para clasificaci\'on es $\sqrt{p} \approx 3{,}6$ para $p=13$
            variables, por lo que el rango $\{2, 4, 6, 8, 10\}$ cubre valores
            por debajo y por encima de dicha referencia \cite{Breiman2001}.
            & $\{2,\, 4,\, 6,\, 8,\, 10\}$ \\
        \bottomrule
        \end{tabularx}
    \end{table}

LDA y la Regresi\'on Log\'istica Multinomial no requieren ajuste expl\'icito
de hiperpar\'ametros en su configuraci\'on est\'andar, por lo que se
emplean directamente como referencias param\'etricas de menor complejidad.
% ─────────────────────────────────────────────────────────────────────────────
\section{Resultados y análisis comparativo}

Tras la fase de entrenamiento y ajuste, los modelos fueron evaluados sobre el conjunto de prueba aislado. A continuación, se presenta la tabla comparativa con los tres mejores algoritmos, ordenados por su rendimiento en la métrica principal.

\vspace{0.5cm}
\begin{table}[h]
    \centering
    \caption{Comparativa global de modelos sobre el conjunto de prueba.}
    \vspace{0.2cm}
    \begin{tabular}{l c c c c c}
        \hline
        \textbf{Modelo} & \textbf{F1-Score} & \textbf{Balanced Acc.} & \textbf{Accuracy} & \textbf{Kappa} & \textbf{MCC} \\
        \hline
        Naive Bayes & 0.9804 & 0.9877 & 0.9808 & 0.9709 & 0.9714 \\
        Random Forest & 0.9804 & 0.9877 & 0.9808 & 0.9709 & 0.9714 \\
        LDA & 0.9623 & 0.9750 & 0.9615 & 0.9419 & 0.9435 \\
        Multinomial & 0.9611 & 0.9754 & 0.9615 & 0.9420 & 0.9441 \\
        Árbol & 0.9444 & 0.9573 & 0.9423 & 0.9122 & 0.9138 \\
        SVM radial & 0.9421 & 0.9630 & 0.9423 & 0.9133 & 0.9179 \\
        kNN & 0.9231 & 0.9503 & 0.9231 & 0.8846 & 0.8911 \\
        Aleatorio - Baseline & 0.4290 & 0.5732 & 0.4231 & 0.1329 & 0.1335 \\
        \hline
    \end{tabular}
    \label{tab:comparativa_modelos}
\end{table}
\vspace{0.5cm}

Como se observa en la Tabla~\ref{tab:comparativa_modelos}, los resultados avalan la viabilidad técnica del sistema. La comparativa de modelos muestra un rendimiento muy elevado en la mayoría de algoritmos, destacando Naive Bayes y Random Forest como los modelos con mejores resultados, alcanzando métricas prácticamente idénticas. Esto sugiere que el problema presenta una estructura bien definida en el espacio de características, donde tanto enfoques probabilísticos como modelos basados en ensamblado son capaces de capturar eficazmente las relaciones entre variables.

LDA y la regresión multinomial presentan un rendimiento muy cercano, lo que refuerza la idea de que una parte importante de la separación entre clases puede modelarse de forma aproximadamente lineal. Por su parte, modelos como SVM radial y kNN no aportan mejoras significativas, lo que indica que la complejidad adicional no resulta determinante en este caso.

El árbol de decisión obtiene un rendimiento inferior, coherente con su mayor varianza en datasets pequeños. Finalmente, el baseline aleatorio se sitúa claramente por debajo del resto, confirmando que los modelos entrenados están capturando una estructura real en los datos.

\begin{table}[h]
    \centering
    \caption{Hiperparámetros óptimos seleccionados mediante Grid Search.}
    \vspace{0.2cm}
    \begin{tabular}{l l}
        \hline
        \textbf{Modelo} & \textbf{Hiperparámetros óptimos} \\
        \hline
        Naive Bayes & \{fL: 0, usekernel: FALSE, adjust: 1\} \\
        Árbol & \{cp: 0.07\} \\
        Random Forest & \{mtry: 2\} \\
        kNN & \{kmax: 10, distance: 2, kernel: optimal\} \\
        SVM radial & \{sigma: 0.01, C: 10\} \\
        \hline
    \end{tabular}
    \label{tab:hiperparametros}
\end{table}
\vspace{0.5cm}

En la Tabla~\ref{tab:hiperparametros} hiperparámetros seleccionados muestran configuraciones relativamente simples, lo que refuerza la idea de que el problema no requiere modelos altamente complejos para alcanzar un buen rendimiento. En particular, destaca el bajo valor de $mtry$ en Random Forest y el valor reducido de $\sigma$ en SVM, indicando fronteras de decisión relativamente suaves y bien definidas.

% ─────────────────────────────────────────────────────────────────────────────
\section{Análisis de errores}

El análisis de errores revela que las observaciones mal clasificadas se concentran principalmente en la clase 2, que en varios casos es predicha como clase 3. Este patrón sugiere la existencia de una región de solapamiento entre ambas clases, donde la separación no es completamente nítida.

\vspace{0.3cm}

En particular, la observación 62 destaca como un caso crítico, ya que es sistemáticamente mal clasificada por todos los modelos. Este comportamiento indica que no se trata de una limitación específica de un algoritmo, sino de una ambigüedad intrínseca en los datos.

\vspace{0.3cm}

\begin{table}[h]
\centering
\caption{Observaciones mal clasificadas comunes entre modelos}
\begin{tabular}{c c c}
\hline
\textbf{Obs} & \textbf{Clase Real} & \textbf{Predicción} \\
\hline
62 & 2 & 3 \\
69 & 2 & 3 \\
119 & 2 & 3 \\
\hline
\end{tabular}
\end{table}

\vspace{0.3cm}

La comparación de la observación 62 con las medias de las clases 2 y 3 confirma que presenta un comportamiento mixto, con ciertas variables más próximas a la clase 3, lo que explica su confusión. En este sentido, se trata de un punto situado en la frontera de decisión, más que de un error aleatorio.

\vspace{0.3cm}

Por otro lado, los outliers detectados previamente mediante PCA no parecen estar directamente relacionados con los errores, ya que únicamente uno de ellos aparece en el conjunto de test y además se clasifica correctamente. No obstante, este aspecto debería analizarse con mayor profundidad en escenarios más amplios, dado que su impacto podría manifestarse en otras particiones o conjuntos de datos.


% ─────────────────────────────────────────────────────────────────────────────
\section{Explicabilidad y validación técnica}
El análisis de importancia de variables mediante Random Forest permite identificar qué características fisicoquímicas tienen mayor influencia en la clasificación. Como se observa, las variables más relevantes son la prolina, la intensidad del color, los flavonoides y el ratio OD280/OD315, seguidas por el alcohol y los compuestos fenólicos totales.

\vspace{0.3cm}

\begin{table}[h]
\centering
\caption{Variables más relevantes según Random Forest}
\begin{tabular}{l c}
\hline
\textbf{Variable} & \textbf{Importancia} \\
\hline
Proline & 100.00 \\
Color_Intensity & 95.01 \\
OD280_OD315 & 84.07 \\
Flavanoids & 75.11 \\
Alcohol & 72.78 \\
\hline
\end{tabular}
\end{table}

\vspace{0.3cm}

Estas variables están directamente relacionadas con la composición fenólica y las características organolépticas del vino, lo que resulta coherente con el dominio enológico. En particular, la prolina y la intensidad del color son conocidos indicadores de diferenciación entre variedades, mientras que los flavonoides y el ratio OD280/OD315 reflejan la concentración de compuestos fenólicos, altamente discriminativos.

\vspace{0.3cm}

En contraste, variables como el \textit{ash} o los \textit{nonflavanoid phenols} presentan una contribución prácticamente nula, lo que indica que aportan poca información relevante para la tarea de clasificación.

\vspace{0.3cm}

Este patrón es consistente con el análisis PCA realizado previamente, donde estas mismas variables mostraban una elevada contribución a las primeras componentes principales. Aunque PCA y Random Forest responden a criterios distintos (varianza explicada frente a capacidad predictiva), la coincidencia observada refuerza la validez del modelo y su interpretación desde el punto de vista enológico.